\section{Conclusion}

In this paper we introduce PersonalAI 2.0 (PAI-2), a novel framework integrating large language model capabilities with graph-based external memory for efficient knowledge retrieval and reasoning. Building upon GraphRAG approach, PAI-2 addresses critical limitations associated with traditional methods, such as inefficiencies in traversing complex knowledge graphs and deficiencies in retrieving precise, context-specific information.

Through a systematic decomposition of queries and dynamic planning of subgraph traversal/retrieval, PAI-2 demonstrates significant improvements over existing methods: LightRAG, RAPTOR, and HippoRAG 2. Evaluations conducted across six benchmarks (Natural Questions, TriviaQA, HotpotQA, 2WikiMultihopQA, MuSiQue, and DiaASQ) highlighted its effectiveness, achieving an average 4\% by LLM-as-a-Judge improvement on 4 out of 6. According to one of the ablation studies, enabled search plan enhancement mechanism give $18\%$ boost (compared to disabled plan enhancement), while advanced graph traversal algorithms give $6\%$ (compared to flatten retriever) boost in retrieval precision.

Additionally, experiments revealed that PAI-2's memory construction algorithm exhibited greater stability than competing methods (KGGen and Wikontic) within 7-14B LLM tier settings, yielding fewer parsing errors and resulted in SOTA on MINE-1 benchmark with $89\%$ information-retention score.

Overall, PAI-2 represents a substantial step forward in combining the expressive power of large language models with the structured data representation offered by knowledge graphs. It expands the way for developing next-generation intelligent agents, capable of delivering both nuanced responses and reliable factual outputs. 

% предложен алгоритм поиска инофрмации на графе который превосходит аналоги на 4 из 6 датасетов
% показано, что при отсутствии стадии перегенрации плана поиска деградирует качество
% выполнен анализ логов и выявлены проблемные места
% выполнено сравгение с предыдущей версией и показан прирост эффективноти при интеграции мехнизма планирования
% проведён дополнительный ablation для исследования надёжности предложенного решения

% TODO